\section{Hiện thực hệ thống}

Trên cơ sở kiến trúc và các quyết định thiết kế ở Chương 3, chương này trình bày quá trình hiện thực hệ thống bằng những thành phần phần mềm và phần cứng cụ thể. Nội dung lần lượt mô tả kiến trúc triển khai, RobotApp, CaregiverApp, quy trình chuẩn bị mô hình AI, firmware của phần thân robot và cách các thành phần phối hợp với nhau trong quá trình vận hành.

\subsection{Kiến trúc hiện thực và ánh xạ công nghệ}

Hình \ref{fig:implementation_architecture} thể hiện các thành phần và công nghệ được sử dụng theo thiết kế ở Chương 3. CaregiverApp cung cấp giao diện cho người chăm sóc. Firebase Authentication, Firebase Realtime Database, dịch vụ LLM và hạ tầng STUN/TURN hỗ trợ xác thực, đồng bộ và liên lạc. RobotApp đảm nhiệm xử lý hình ảnh, hội thoại và điều phối tại robot. Phần thân gồm ESP32, cảm biến góc, mạch công suất và hai động cơ BLDC; firmware sử dụng BLE GATT để nhận lệnh và SimpleFOC để điều khiển động cơ.

Firebase RTDB lưu dữ liệu dưới dạng cây JSON và đồng bộ trạng thái chức năng, lệnh điều khiển cùng thông tin thiết lập cuộc gọi gần thời gian thực. Âm thanh và hình ảnh được truyền bằng WebRTC, trực tiếp giữa hai thiết bị hoặc qua máy chủ TURN. Dịch vụ Groq xử lý hội thoại và các yêu cầu phân loại ngôn ngữ. RobotApp tiếp nhận kết quả AI, phối hợp các trạng thái xử lý và chuyển lệnh chuyển động từ RTDB hoặc chức năng bám theo đến ESP32 qua BLE. Firmware thực hiện lệnh, điều khiển động cơ và dừng khi không còn nhận được cập nhật.

\begin{figure}[H]
    \centering
    \includegraphics[width=\textwidth]{System_Implementation_Architecture_Overview_Style_Figma.png}
    \caption{Kiến trúc hiện thực của hệ thống}
    \label{fig:implementation_architecture}
\end{figure}

\subsection{Môi trường phát triển và tổ chức các thành phần}

Hệ thống gồm bốn thành phần được phát triển độc lập nhưng sử dụng cấu trúc dữ liệu và quy ước trao đổi thống nhất. RobotApp và CaregiverApp được xây dựng bằng Kotlin, giao diện XML và ViewBinding. RobotApp sử dụng package \texttt{com.example.datn\_v1}, còn CaregiverApp sử dụng package \texttt{com.example.datn\_caregiver}; cả hai cùng đặt \texttt{minSdk=23}, \texttt{targetSdk=36} và Java/Kotlin JVM target 11. Firmware được xây dựng dưới dạng chương trình Arduino cho ESP32. Thành phần AI gồm mô hình sau huấn luyện, mô hình ONNX trung gian, mô hình TFLite và các chương trình chuyển đổi, kiểm tra suy luận.

\begin{table}[htbp]
\centering
\caption{Các thành phần hiện thực của hệ thống}
\label{tab:source_blocks}
\small
\begin{tabularx}{\textwidth}{|p{3.1cm}|p{5.1cm}|X|}
\hline
\textbf{Khối} & \textbf{Công nghệ chính} & \textbf{Vai trò} \\
\hline
RobotApp & Android/Kotlin, CameraX, TFLite, OpenCV, Firebase, WebRTC, BLE, OkHttp & Xử lý AI, hội thoại, cuộc gọi và kết nối với thân robot \\
\hline
CaregiverApp & Android/Kotlin, Firebase Auth/RTDB, WebRTC & Ứng dụng cho người thân quản lý, giám sát và điều khiển \\
\hline
Mô hình AI & Python, Keras/TensorFlow, Ultralytics, ONNX, onnx2tf, TFLite & Chuẩn bị dữ liệu, huấn luyện và chuyển đổi mô hình cho Android \\
\hline
Firmware & ESP32, NimBLE, SimpleFOC, AS5600, 3PWM & Tiếp nhận, phân tích lệnh qua BLE và điều khiển hai động cơ BLDC \\
\hline
\end{tabularx}
\end{table}

RobotApp sử dụng CameraX, các gói TensorFlow Lite (CPU, GPU, Support và Select TF Ops), OpenCV Android, Firebase Authentication, Firebase Database, WebRTC native, OkHttp và Kotlin Coroutines. CaregiverApp sử dụng Firebase Authentication, Firebase Database, Android Credentials, RecyclerView và WebRTC. Tệp mô hình TFLite được cấu hình \texttt{noCompress} để giữ nguyên dạng khi đóng gói, cho phép truy cập bằng cơ chế ánh xạ tệp vào bộ nhớ.

\subsection{Hiện thực RobotApp}

\subsubsection{Khởi động, điều hướng và vòng đời dịch vụ}

\texttt{LoginActivity} là màn hình khởi động của RobotApp. Sau khi xác thực, \texttt{MainActivity} quản lý các fragment hồ sơ, theo dõi, chatbot và cuộc gọi, đồng thời lắng nghe nhánh \texttt{webrtc\_signal} để hiển thị yêu cầu gọi đến. \texttt{ReminderService} theo dõi lịch nhắc nhở; \texttt{TrackingService} duy trì thông báo của dịch vụ chạy nền và chuyển sự kiện té ngã từ khối theo dõi đến chatbot. Việc thu nhận hình ảnh và suy luận AI được thực hiện trong \texttt{tracking\_fragment}. Tệp AndroidManifest khai báo các quyền sử dụng camera, microphone, Internet, thông báo, dịch vụ nền và Bluetooth theo phiên bản hệ điều hành.

Quyền Bluetooth được khai báo theo từng nhóm phiên bản hệ điều hành. Android 11 trở xuống sử dụng \texttt{BLUETOOTH}, \path|BLUETOOTH_ADMIN| và quyền vị trí để quét thiết bị; từ Android 12, ứng dụng sử dụng \path|BLUETOOTH_SCAN| và \path|BLUETOOTH_CONNECT| theo hướng dẫn của Android \reportcite{ref:android_permissions}. Camera được khai báo \texttt{required=false} để không ngăn cài đặt chỉ vì thiết bị thiếu cờ tính năng tương ứng, nhưng các chức năng thị giác chỉ hoạt động khi có camera.

\subsubsection{TrackingFragment và quản lý camera}

\texttt{tracking\_fragment.kt} điều phối CameraX, mô hình ước lượng tư thế, bộ theo dõi, LSTM, bộ hậu kiểm và chức năng bám người. CameraX cung cấp \texttt{Preview} để hiển thị hình ảnh xem trước và \texttt{ImageAnalysis} để phân tích khung hình. Bộ phân tích chuyển dữ liệu từ \texttt{ImageProxy} sang ảnh bitmap và tensor, sau đó giải phóng khung hình đã xử lý. Hai chức năng camera được liên kết với vòng đời của fragment theo cơ chế CameraX \reportcite{ref:android_camera}.

Khi không có cuộc gọi, CameraX được gắn với vòng đời của fragment. Trước khi WebRTC khởi động, CameraX được tháo liên kết để tránh hai thành phần cùng mở camera. \texttt{video\_call\_fragment} tiếp quản camera và gửi một bản sao \texttt{VideoFrame} cục bộ về khối theo dõi. Khung hình được chuyển từ I420 sang RGB, giảm kích thước và xử lý trên một executor riêng; nếu lượt suy luận trước chưa kết thúc, khung hình mới bị bỏ qua thay vì tiếp tục xếp hàng. Khi cuộc gọi kết thúc, bộ thu hình của WebRTC được dừng và giải phóng trước khi CameraX được liên kết trở lại.

Bên cạnh việc quản lý nguồn hình ảnh, \texttt{tracking\_fragment} duy trì một đối tượng \texttt{RobotControlManager} để điều khiển kết nối với ESP32. Chức năng gọi video sử dụng chung bộ quản lý này, giúp tập trung việc kiểm tra lệnh và xử lý dừng robot.

\subsubsection{Nạp và suy luận mô hình YOLO26n-Pose}

Các quy ước về bounding box, keypoint và hậu xử lý ở Mục \ref{sec:yolo_pose_output_representation} và Mục \ref{sec:yolo_pose_inference_pipeline} được hiện thực trong \texttt{tracking\_fragment.kt}. Chuỗi xử lý bao gồm lựa chọn mô hình, khởi tạo bộ suy luận, tiền xử lý khung hình, giải mã tensor và chuyển kết quả đến SORT, LSTM cùng lớp hiển thị.

\paragraph{Mô hình triển khai và lựa chọn bộ tăng tốc.}
RobotApp đóng gói hai biến thể mô hình trong thư mục \texttt{assets}. Biến thể đầu ra thô có kích thước tensor tĩnh được ưu tiên để khai thác GPU; biến thể đầu cuối được sử dụng khi mô hình thứ nhất không thể khởi tạo. Bảng \ref{tab:yolo_runtime_artifacts} trình bày quy ước đầu ra mà bộ giải mã hỗ trợ.

\begin{table}[htbp]
\centering
\caption{Hai mô hình YOLO26n-Pose được đóng gói trong RobotApp}
\label{tab:yolo_runtime_artifacts}
\renewcommand{\arraystretch}{1.12}
\footnotesize
\begin{tabularx}{\textwidth}{|>{\raggedright\arraybackslash}p{3.55cm}|p{1.7cm}|p{2.45cm}|X|}
\hline
\textbf{Tệp mô hình} & \textbf{Kích thước} & \textbf{Đầu ra} & \textbf{Vai trò} \\
\hline
\path|yolo26n-pose_320_gpu.tflite| & 11.838.836 byte & \texttt{[1,56,2100]} & Đầu ra thô, tensor tĩnh; được ưu tiên để suy luận bằng GPU \\
\hline
\path|yolo26n-pose_320.tflite| & 12.031.560 byte & \texttt{[1,300,57]} & Đầu ra E2E; mô hình dự phòng với số lượng kết quả giới hạn \\
\hline
\end{tabularx}
\end{table}

Sau khi chọn tệp, \texttt{setupAIModel()} lần lượt thử bộ tăng tốc GPU, NNAPI và cuối cùng là CPU với bốn luồng. Nếu GPU không khởi tạo được, tài nguyên của bộ tăng tốc này được giải phóng trước khi chuyển sang NNAPI. Thay vì cố định kích thước theo tên tệp, chương trình đọc \texttt{getInputTensor(0).shape()} và \texttt{getOutputTensor(0).shape()}, sau đó cập nhật kích thước đầu vào và cấp phát bộ đệm đầu ra theo hình dạng tensor thực tế. Nhờ vậy, cùng một chuỗi xử lý có thể tiếp nhận các mô hình được xuất theo những quy ước khác nhau. TFLite được lựa chọn vì phù hợp với việc triển khai mô hình trên thiết bị di động và thiết bị biên \reportcite{ref:tflite}.

\paragraph{Tiền xử lý khung hình.}
CameraX sử dụng \texttt{STRATEGY\_KEEP\_ONLY\_LATEST} để ưu tiên khung hình mới nhất. Bộ điều phối suy luận bỏ qua khung hình khi chưa đủ khoảng thời gian cấu hình hoặc lượt xử lý trước chưa hoàn tất, nhờ đó hạn chế tích lũy ảnh cũ trong hàng đợi. Trong cuộc gọi, bản sao \texttt{VideoFrame} của WebRTC được đưa đến cùng bộ thực thi tác vụ và hàm suy luận, cho phép hai nguồn hình ảnh sử dụng chung quy trình xử lý.

Ảnh bitmap được nạp vào \texttt{TensorImage} kiểu \texttt{float32}, đổi trực tiếp về kích thước $320\times320$ bằng \texttt{ResizeOp.ResizeMethod.BILINEAR}, rồi chuẩn hóa giá trị RGB từ $[0,255]$ về $[0,1]$ bằng \texttt{NormalizeOp(0f,255f)}. Ứng dụng không thêm phần đệm ảnh theo phương pháp letterbox. Với khung hình camera $W_c\times H_c$ và $S=320$, ánh xạ từ tọa độ ảnh camera sang ảnh đầu vào mô hình là
\begin{equation}
x_i=x_c\frac{S}{W_c},\qquad y_i=y_c\frac{S}{H_c},
\end{equation}
còn phép đổi ngược là $x_c=x_iW_c/S$, $y_c=y_iH_c/S$. Hai hệ số theo hai trục được giữ riêng vì khung hình nguồn không nhất thiết là ảnh vuông.

\paragraph{Nhận dạng và giải mã cấu trúc tensor.}
Đầu ra thô chứa dự đoán tại ba mức bước lấy mẫu (stride) 8, 16 và 32. Với ảnh $320\times320$, tổng số vị trí dự đoán là
\begin{equation}
N=\left(\frac{320}{8}\right)^2+
  \left(\frac{320}{16}\right)^2+
  \left(\frac{320}{32}\right)^2=2100.
\end{equation}
Vì mô hình chỉ nhận diện lớp người, mỗi vị trí dự đoán gồm 4 giá trị hộp, 1 độ tin cậy và $17\times3$ giá trị keypoint, tổng cộng 56 kênh. Mô hình đầu ra thô trả tensor \texttt{[1,56,2100]}; hàm \texttt{getOutputVal()} xử lý sự khác biệt giữa cách sắp xếp này và dạng \texttt{[1,2100,56]}, nhờ đó bộ giải mã luôn truy cập theo cặp chỉ số dự đoán--kênh. Bốn kênh đầu biểu diễn $(c_x,c_y,w_b,h_b)$, kênh 4 là độ tin cậy, còn 51 giá trị từ kênh 5 tạo thành 17 bộ ba $(x_j,y_j,p_j)$.

Đầu ra đầu cuối (E2E) được nhận biết khi hai chiều còn lại sau chiều kích thước lô có giá trị $300$ và $57$. Mỗi hàng gồm bốn trường đầu biểu diễn $(x_{\min},y_{\min},x_{\max},y_{\max})$, trường 4 là độ tin cậy, trường 5 là nhãn lớp và các trường từ 6 chứa dữ liệu của 17 điểm khớp. Hai dạng đầu ra cần được giải mã riêng: đọc đầu ra thô theo quy ước E2E sẽ làm tọa độ tâm và kích thước bounding box bị hiểu thành tọa độ hai góc, đồng thời làm lệch vị trí các kênh điểm khớp.

Hai mô hình có thể trả tọa độ chuẩn hóa hoặc tọa độ pixel trong không gian đầu vào. Hàm \texttt{detectCoordSystem()} khảo sát kênh tọa độ ở lần suy luận đầu tiên: nếu giá trị dương lớn nhất không vượt quá 2, tọa độ được xem là đã chuẩn hóa; ngược lại, kết quả được xem là tọa độ pixel của ảnh mô hình. Trường hợp thứ nhất được nhân trực tiếp với $W_c,H_c$, còn trường hợp thứ hai được nhân với $W_c/320,H_c/320$. Quy tắc nhận biết này đáp ứng hai mô hình được sử dụng trong đề tài; khi bổ sung quy trình xuất mô hình mới, hệ tọa độ nên được khai báo rõ bằng siêu dữ liệu để tránh suy đoán.

Đầu ra thô được đổi từ dạng $cxcywh$ sang $xyxy$, sau đó loại các dự đoán có độ tin cậy dưới $0{,}20$. Hàm \texttt{nms()} sắp xếp kết quả theo điểm tin cậy và loại hộp có IoU lớn hơn $0{,}45$. Kết quả E2E cũng đi qua hàm này như một lớp lọc bổ sung, mặc dù mô hình đầu cuối đã thực hiện bước lựa chọn dự đoán. Mỗi keypoint được giữ kèm độ tin cậy và chỉ được đánh dấu \texttt{visible} khi $p_j>0{,}30$.

\paragraph{Chuyển đổi tọa độ và hiển thị kết quả.}
Quy trình xử lý sử dụng bốn hệ tọa độ như trình bày trong Bảng \ref{tab:yolo_coordinate_systems}. Bounding box và các điểm khớp phải được chuyển đổi nhất quán giữa các hệ này. Các phép tính theo dõi và điều khiển sử dụng tọa độ ảnh camera; tọa độ trên giao diện chỉ phục vụ hiển thị.

\begin{table}[htbp]
\centering
\caption{Các hệ tọa độ trong quy trình xử lý YOLO26n-Pose của RobotApp}
\label{tab:yolo_coordinate_systems}
\renewcommand{\arraystretch}{1.12}
\small
\begin{tabularx}{\textwidth}{|p{2.35cm}|p{2.35cm}|X|X|}
\hline
\textbf{Không gian} & \textbf{Miền tọa độ} & \textbf{Vai trò} & \textbf{Dữ liệu sử dụng} \\
\hline
Ảnh camera $\mathcal{C}$ & $[0,W_c]\times[0,H_c]$ pixel & Khung hình gốc nhận từ CameraX hoặc WebRTC & SORT, LSTM và điều khiển bám người \\
\hline
Ảnh mô hình $\mathcal{I}$ & $[0,320]\times[0,320]$ pixel & Ảnh vuông sau \texttt{ResizeOp} & Suy luận và đầu ra dạng pixel \\
\hline
Chuẩn hóa $\mathcal{N}$ & $[0,1]\times[0,1]$ & Biểu diễn không phụ thuộc độ phân giải & Đầu ra chuẩn hóa và các tỉ lệ hình học \\
\hline
Màn hình $\mathcal{V}$ & $[0,W_v]\times[0,H_v]$ pixel & \texttt{BoundingBoxOverlay} chồng lên preview & Vẽ hộp, skeleton, nhãn và track ID \\
\hline
\end{tabularx}
\end{table}

Camera trước được hiển thị như gương. Sau khi đổi về pixel camera, tọa độ ngang của keypoint được thay bằng $x'_c=W_c-x_c$. Với bounding box, hai cạnh phải đổi đồng thời:
\begin{equation}
x'_{\min}=W_c-x_{\max},\qquad x'_{\max}=W_c-x_{\min}.
\end{equation}
Nếu chỉ phản chiếu từng cạnh mà không hoán đổi, kết quả sẽ có cạnh phải nhỏ hơn cạnh trái. Tọa độ $y$ không thay đổi.

\texttt{BoundingBoxOverlay.onDraw()} tiếp tục ánh xạ pixel camera sang View theo center-crop giống \texttt{PreviewView}. Với
\begin{equation}
s_v=\max\!\left(\frac{W_v}{W_c},\frac{H_v}{H_c}\right),\quad
\Delta_x=\frac{W_v-s_vW_c}{2},\quad
\Delta_y=\frac{H_v-s_vH_c}{2},
\end{equation}
một điểm được vẽ tại $x_v=s_vx'_c+\Delta_x$, $y_v=s_vy_c+\Delta_y$. Phần bù có thể âm vì center-crop cắt bớt một chiều. Tọa độ View chỉ phục vụ hiển thị; các đại lượng như tâm mục tiêu, mức chiếm khung và khoảng cách keypoint vẫn được tính trong không gian camera rồi chuẩn hóa bởi $W_c,H_c$.

\paragraph{Chuyển kết quả đến các module tiếp theo.}
Sau khi giải mã, mỗi đối tượng gồm bounding box \texttt{bbox}, độ tin cậy \texttt{score}, nhãn lớp \texttt{classId} và 17 điểm khớp trước khi được chuyển đến \texttt{SortTracker}. Kết quả theo dõi được dùng để xác định mục tiêu ID0, tạo vector 69 đặc trưng cho LSTM, tính lệnh bám theo và hiển thị thông tin lên màn hình. Dữ liệu được thống nhất trong hệ tọa độ camera, nên các module phía sau không cần xử lý lại sự khác biệt về dạng đầu ra, thứ tự các chiều tensor hoặc quy ước tọa độ của từng mô hình.

\subsubsection{SORT, Kalman, Hungarian và Re-ID}

Các lớp \texttt{SortTracker}, \texttt{KalmanBoxTracker}, \texttt{HungarianAlgorithm} và \texttt{ReIdModule} hiện thực phương pháp theo dõi dựa trên kết quả phát hiện. Mỗi đối tượng gồm \texttt{bbox}, \texttt{score}, \texttt{classId} và các keypoint. Bộ lọc Kalman dự báo vùng bao ở khung hình kế tiếp, còn thuật toán Hungarian tìm phương án ghép có tổng chi phí nhỏ nhất. Phương pháp này kế thừa cấu trúc gọn của SORT \reportcite{ref:sort}, đồng thời bổ sung tín hiệu tư thế và đặc trưng ngoại hình để phù hợp hơn với bài toán trên điện thoại.

Hàm chi phí ghép đối tượng gồm bốn thành phần: 30\% sai khác IoU, 30\% khoảng cách tâm, 30\% khoảng cách giữa các điểm khớp và 10\% thay đổi diện tích bounding box. Ngưỡng IoU cơ sở của bộ theo dõi là 0,25; đối tượng mới cần hai lần ghép thành công để được xác nhận. Khi không tìm thấy kết quả phát hiện phù hợp, định danh ID0 được duy trì tối đa 45 khung hình, còn các định danh khác được giữ tối đa 10 khung hình.

Khối tái định danh trích xuất biểu đồ phân bố màu (histogram) trong không gian Hue--Saturation--Value (HSV) từ vùng thân người, nhằm giảm ảnh hưởng của hậu cảnh và thay đổi góc mặt. Đặc trưng màu không được cập nhật khi thiếu điểm khớp tin cậy, vì vùng cắt không chính xác có thể chứa hình ảnh của người khác. Khi một kết quả phát hiện phù hợp với đặc trưng đã lưu, \texttt{restoreWithId} khôi phục định danh ban đầu. Giải pháp này có chi phí tính toán thấp hơn việc bổ sung mạng trích xuất đặc trưng ngoại hình như Deep SORT \reportcite{ref:deepsort}, nhưng dễ bị ảnh hưởng khi nhiều người mặc trang phục giống nhau hoặc ánh sáng thay đổi.

\subsubsection{Duy trì mục tiêu ID0 và tạo lệnh bám theo}

Người mục tiêu được gán định danh \texttt{PRIMARY\_TRACK\_ID=0}. Việc lựa chọn ban đầu dựa trên diện tích bounding box, độ tin cậy phát hiện, tỷ lệ điểm khớp quan sát được, khoảng cách đến tâm ảnh và mức độ bị cắt ở biên. Định danh ID0 được dùng chung cho phân tích té ngã và điều khiển bám theo. Khi mất dấu, hệ thống tìm lại mục tiêu trong khoảng 2 s; đối tượng được xét phải cách vị trí cuối không quá 35\% đường chéo ảnh, nằm trong vùng trung tâm và được quan sát ổn định qua ba khung hình. Trong quá trình xử lý sự kiện nghi té ngã, hệ thống không tự chuyển sang một người khác.

Chức năng bám theo ước lượng độ gần từ kích thước bounding box và khoảng cách giữa các mốc vai, hông, gối, mắt cá chân. Giá trị sau khi làm mượt được chia thành ba vùng quá xa, phù hợp và quá gần. Ngưỡng đi vào và rời khỏi mỗi vùng khác nhau để tránh chuyển liên tục giữa tiến, dừng và lùi khi giá trị dao động gần biên. \texttt{TrackingMotionPolicy} điều chỉnh thành phần tịnh tiến theo độ gần: mức tiến tăng từ 0,60 đến 0,90, mức lùi khi mục tiêu quá gần thay đổi từ 0,45 đến 0,60, còn mức lùi để quan sát lại phần chân là 0,30. Đây là các độ lớn của lệnh điều khiển chuẩn hóa, không phải vận tốc đo bằng m/s. Khi robot vừa quay vừa tiến hoặc chưa quan sát ổn định mắt cá chân, mức tiến được giảm để hạn chế mất dấu. Sai lệch ngang giữa tâm bounding box và tâm ảnh tạo thành phần quay liên tục, có vùng chết và giới hạn biên để giảm dao động.

Một nhánh hình học riêng phân loại tư thế phục vụ điều khiển, không thay đổi kết quả FALL/NO FALL của LSTM. Nhánh này kết hợp tỷ lệ vùng bao, hướng phân bố keypoint, quan hệ vai--hông và bằng chứng từ thân dưới; \texttt{PostureGeometry} bổ sung trường hợp người nằm dọc theo hướng nhìn của camera nên hình chiếu có thể trông gần giống tư thế đứng. Khi tư thế nằm đã ổn định hoặc một sự cố té ngã đang được xử lý, robot không tiếp tục tiến mà chỉ dừng, lùi chậm hoặc quay tại chỗ để đưa người vào vùng quan sát phù hợp. Mọi vector điều khiển cuối cùng vẫn được gửi qua \texttt{RobotControlManager} để kiểm tra và chuyển thành vận tốc hai bánh.

\subsubsection{LSTM phát hiện té ngã}

\texttt{LstmFallDetector} nạp mô hình \texttt{best\_model\_lstm\_v2.tflite} có dung lượng khoảng 63,9 KB và sử dụng hai luồng CPU. Dữ liệu của mỗi đối tượng được lưu trong \texttt{FallDetectionQueue}; khi tích lũy đủ 15 khung hình, hàng đợi tạo tensor đầu vào $1\times15\times69$ kiểu số thực. Tensor đầu ra $1\times2$ biểu diễn xác suất của hai lớp hoạt động bình thường và té ngã. So với tệp Keras gốc khoảng 193,8 KB, tệp TFLite có dung lượng nhỏ hơn khi đóng gói vào ứng dụng.

Kết quả được gán nhãn FALL khi xác suất FALL lớn hơn xác suất NO FALL, sau đó chuyển đến \texttt{FallVerificationEngine} để kiểm tra bổ sung. Chỉ khi có năm kết quả FALL liên tiếp, hệ thống mới ghi nhận một chuỗi nghi té ngã. Khi bộ theo dõi khôi phục thành công ID0, hàng đợi đặc trưng vẫn được gắn với người mục tiêu. Khả năng phân tích quan hệ thời gian của LSTM \reportcite{ref:lstm} được kết hợp với bước hậu kiểm trước khi tạo cảnh báo.

\subsubsection{Hậu kiểm té ngã}

\texttt{FallVerificationEngine} tiếp nhận \texttt{FrameEvidence}, gồm trạng thái quan sát mục tiêu, vị trí tương đối của thân trên, số điểm khớp, mức độ bounding box nằm trọn trong ảnh và tỷ lệ hình dạng. Sau khi chuỗi FALL kết thúc, bộ hậu kiểm xem xét vị trí mũi, vai ở các khung hình tiếp theo. Ba khung hình liên tiếp có các điểm thân trên trong vùng 70\% phía trên ảnh sẽ loại bỏ tình huống nghi té ngã; ba khung hình liên tiếp có các điểm này nằm dưới ngưỡng, hoặc 12 khung hình liên tiếp thiếu dữ liệu quan sát, sẽ dẫn đến bước hỏi xác nhận. Khi sự kiện đã được loại bỏ hoặc xác nhận, bộ hậu kiểm chờ 10 kết quả NO FALL liên tiếp trước khi trở lại trạng thái giám sát.

Bộ hậu kiểm trả về một trong ba kết quả: chưa đủ dữ liệu để kết luận, loại bỏ tình huống nghi té ngã qua \texttt{Rejected(reason)}, hoặc yêu cầu hỏi xác nhận qua \texttt{AskForConfirmation(reason, score, evidence)}. Việc tách phần xử lý khỏi giao diện cho phép kiểm thử riêng từng trường hợp bằng chuỗi dữ liệu đầu vào mà không cần camera. Trạng thái \texttt{isIncidentActive} cho biết sự kiện đang được xử lý, còn \texttt{canAcquireNewTarget} quy định khi nào được phép chọn mục tiêu mới. Các hàm \texttt{resolveAsSafe}, \texttt{resolveAsDanger} và \texttt{reset} cập nhật kết quả xác nhận hoặc đặt lại bộ xử lý.

\begin{figure}[htbp]
    \centering
    \includegraphics[width=0.88\textwidth]{1.jpeg}
    \caption{Ví dụ khung hình té ngã trong dữ liệu tham khảo}
    \label{fig:fall_sample}
\end{figure}

\subsubsection{Luồng xác nhận và cảnh báo}

Trong mỗi sự cố, TrackingService chỉ chuyển một \texttt{FallEvent} đến chatbot. Robot dừng chuyển động, phát câu hỏi bằng TTS và lắng nghe câu trả lời trong tối đa 10 s. Những cụm từ thể hiện rõ trạng thái an toàn hoặc nguy hiểm được so khớp tại thiết bị để giảm độ trễ và hạn chế phụ thuộc vào LLM. Với câu trả lời không rõ nghĩa, hệ thống gọi mô hình phân loại theo chế độ không truyền luồng và yêu cầu kết quả có cấu trúc. Nếu người dùng xác nhận an toàn, sự cố được kết thúc tại robot; các trường hợp nguy hiểm, lỗi phân loại hoặc không có phản hồi sẽ được ghi vào \texttt{fall\_alerts}, sau đó cuộc gọi đến người chăm sóc được thực hiện khi TTS phát xong.

Bản ghi cảnh báo lưu thời điểm, trạng thái xử lý, phản hồi của người dùng và các thông tin kiểm tra bổ sung như \path|verificationReason|, \path|visibleKeypoints|, \path|bboxFullyInside| và \path|upperBodyYRatio|. Các trường này lần lượt mô tả lý do yêu cầu xác nhận, số điểm khớp quan sát được, bounding box có nằm hoàn toàn trong ảnh hay không và vị trí tương đối của phần thân trên. Chúng hỗ trợ xem lại quá trình xử lý sự kiện, nhưng chưa giải thích đầy đủ cách mô hình LSTM đưa ra dự đoán.

\subsubsection{Chatbot giọng nói}

\path|chatbot_fragment| phối hợp ba thành phần: \path|SpeechRecognizerManager| nhận dạng tiếng nói, \path|GroqChatService| xử lý hội thoại và \path|AndroidTtsService| tổng hợp tiếng nói. Giao diện hiển thị trạng thái hoạt động, biểu tượng lớn, tên robot và nội dung đang được nhận dạng hoặc phát ra. Hồ sơ tại \texttt{Users/\{uid\}} cung cấp cách xưng hô, tên người chăm sóc, tình trạng và sở thích để bổ sung ngữ cảnh cho câu trả lời.

\texttt{GroqChatService} gọi điểm cuối \texttt{/openai/v1/chat/completions} qua OkHttp. Hội thoại dạng truyền luồng sử dụng mô hình \texttt{qwen/qwen3.6-27b}, trong khi các tác vụ phân loại ngắn sử dụng \texttt{openai/gpt-oss-20b}. Lịch sử được giới hạn ở 20 thông điệp để kiểm soát độ dài ngữ cảnh. Với phản hồi theo cơ chế Server-Sent Events (SSE) có tiền tố \texttt{data:}, dịch vụ đọc nội dung tại \texttt{choices[0].delta.content}, phát từng phần qua Kotlin Flow và lưu câu trả lời hoàn chỉnh khi luồng kết thúc.

Các tác vụ liên quan đến an toàn không phụ thuộc hoàn toàn vào văn bản do LLM sinh ra. Hàm phân loại phản hồi té ngã kiểm tra trước tập từ khóa an toàn và nguy hiểm; API chỉ được gọi khi câu trả lời không rõ nghĩa, sau đó kết quả JSON được phân tích theo cấu trúc định trước. Hàm \texttt{detectCallIntent} áp dụng cách tương tự khi kết hợp từ khóa, tên hoặc quan hệ của người chăm sóc và mô hình phân loại dự phòng. Khóa API được truyền vào lớp dịch vụ và cần được cung cấp từ cấu hình bên ngoài kho mã hoặc thông qua máy chủ trung gian.

\subsubsection{Nhắc nhở theo lịch}

\texttt{ReminderService} lắng nghe lịch tại \texttt{Schedules/\{uid\}} và hồ sơ tại \texttt{Users/\{uid\}}. Khi đến thời điểm nhắc, dịch vụ tạo bản ghi dưới \texttt{reminder\_logs/\{uid\}} và gửi sự kiện đến chatbot để phát lời nhắc. Người cao tuổi có thể xác nhận, từ chối hoặc không phản hồi. Bản ghi được cập nhật trạng thái, câu trả lời, số lần nhắc lại và thời điểm xác nhận. Việc lưu riêng từng lần thực hiện cho phép người chăm sóc xem lịch sử mà không thay đổi nội dung lịch gốc.

\subsubsection{WebRTC trong RobotApp}

\texttt{video\_call\_fragment} khởi tạo \texttt{PeerConnectionFactory}, nguồn âm thanh và hình ảnh, thành phần hiển thị video cùng tham chiếu đến nhánh \texttt{webrtc\_signal}. Khi chủ động gọi, RobotApp ghi bản đề nghị \texttt{offer} và các địa chỉ ứng viên ICE vào \texttt{robotCandidates}. Khi nhận cuộc gọi, ứng dụng đọc \texttt{offer}, tạo bản trả lời \texttt{answer} và theo dõi \texttt{caregiverCandidates}. Các bộ lắng nghe được đăng ký theo từng giai đoạn thiết lập phiên và hủy khi không còn sử dụng.

WebRTC đảm nhiệm mã hóa và truyền âm thanh, hình ảnh; Firebase trao đổi SDP, ICE và trạng thái cuộc gọi \reportcite{ref:webrtc}. Trường \texttt{emergencyReason=fall} phân biệt cuộc gọi liên quan đến té ngã với cuộc gọi thông thường. Khi kết thúc, ứng dụng đặt trạng thái thành \texttt{ended}, giải phóng luồng dữ liệu, bộ thu hình, thành phần hiển thị và \texttt{PeerConnection}, sau đó trả quyền sử dụng camera cho chức năng theo dõi.

\subsubsection{RobotControlManager và BLE}

\texttt{RobotControlManager} lắng nghe nhánh \texttt{robot\_control}. Hàm ánh xạ lệnh được tách riêng để kiểm thử bằng bốn điểm chuẩn. Mỗi giá trị đầu vào được kiểm tra tính hữu hạn, giới hạn biên, độ tuổi và độ lệch thời gian về tương lai. Chuỗi lệnh được định dạng với \texttt{Locale.US} để dấu phân cách thập phân luôn là dấu chấm:

\begin{lstlisting}[caption={Định dạng lệnh BLE ở mức logic}]
left  = clamp((steering * 0.15 + translation) * 20, -20, 20)
right = clamp((steering * 0.15 - translation) * 20, -20, 20)
payload = "L%.2fR%.2f\n".format(Locale.US, left, right)
\end{lstlisting}

\texttt{RobotControlManager} tìm thiết bị theo tên \texttt{OhmniRobot} và UUID đã quy định. Thời gian quét tối đa là 10 s; khi cần khôi phục kết nối, ứng dụng thử kết nối lại sau 3 s. Sau khi tìm được dịch vụ và đặc tính nhận lệnh, ứng dụng sử dụng kiểu ghi \texttt{WRITE\_TYPE\_NO\_RESPONSE}, với cách gọi API phù hợp từng phiên bản Android. Nếu không có lệnh mới từ RTDB trong 500 ms, bộ giám sát thời gian nhận lệnh gửi lệnh dừng xuống ESP32. Hàm \texttt{stop()} cũng gửi lệnh dừng trước khi hủy các bộ lắng nghe, hàm xử lý sự kiện, tác vụ quét và kết nối GATT.

\subsection{Hiện thực CaregiverApp}

\subsubsection{Khởi động, xác thực và hồ sơ}

CaregiverApp sử dụng \texttt{LoginActivity} và \texttt{SignUpActivity} cho quá trình xác thực, \texttt{MainActivity} làm màn hình khởi động chính, còn \texttt{ProfileFragment} phục vụ quản lý hồ sơ. Firebase Authentication cung cấp UID; hồ sơ được lưu tại \texttt{Users/\{uid\}} để RobotApp cùng truy cập. \texttt{MainActivity} đồng thời khởi động \texttt{IncomingCallService} nhằm tiếp nhận cuộc gọi khi ứng dụng chạy nền. Do Activity này được khai báo là launcher, bước kiểm tra phiên xác thực ngay khi khởi động cần được hoàn thiện trước khi triển khai cho người dùng thực tế.

\subsubsection{Bật/tắt chức năng theo dõi}

\path|tracking_fragment| cập nhật hai giá trị trạng thái \path|robot_follow_enabled| và \path|ai_fall_detection_enabled| để bật, tắt các chức năng tương ứng. Fragment theo dõi \texttt{.info/connected} để hiển thị tình trạng kết nối RTDB. Khi tắt bám theo tự động, CaregiverApp đồng thời gửi lệnh $x=0,y=0$ kèm thời điểm hiện tại để robot dừng chuyển động.

\subsubsection{Lịch và nhật ký}

Danh sách lịch tại \path|Schedules/uid| được đọc bởi \texttt{ScheduleFragment}. Người chăm sóc tạo hoặc cập nhật lịch qua \texttt{AddScheduleFragment}; \texttt{ScheduleAdapter} hiển thị danh sách bằng RecyclerView và liên kết công tắc bật, tắt với trường \texttt{active}. Nhật ký lời nhắc, cảnh báo té ngã và cuộc gọi có các fragment cùng bộ hiển thị danh sách riêng. Cảnh báo và cuộc gọi được truy vấn theo thời điểm trong bảy ngày gần nhất.

\subsubsection{Tiếp nhận cuộc gọi khi ứng dụng chạy nền}

\texttt{IncomingCallService} là dịch vụ chạy nền có thông báo thường trực (foreground service), lắng nghe nhánh \texttt{webrtc\_signal}. Khi nhận \texttt{status=calling} và \texttt{callerRole=robot}, dịch vụ tạo thông báo và có thể mở \texttt{IncomingCallActivity} toàn màn hình. Nếu có \texttt{emergencyReason=fall}, nội dung thông báo cho biết cuộc gọi được tạo sau tình huống chưa xác nhận an toàn. Tệp Manifest khai báo các quyền liên quan đến thông báo, dịch vụ nền và hiển thị toàn màn hình; khả năng hiển thị phụ thuộc vào quyền được cấp và chính sách của hệ điều hành.

\subsubsection{Gọi video và điều khiển robot từ xa}

\texttt{video\_call\_fragment} xử lý cả vai trò gọi đi và nhận cuộc gọi, hiển thị hình ảnh cục bộ, từ xa và trao đổi thông tin thiết lập phiên. Trong màn hình gọi, cần điều khiển ảo tính $x=\Delta y/r$, $y=-\Delta x/r$, giới hạn giá trị trong $[-1,1]$ và giới hạn tần suất gửi lệnh khoảng 20 Hz. Khi người dùng bắt đầu điều khiển, ứng dụng tắt bám theo tự động để tránh hai nguồn lệnh cùng hoạt động. Các sự kiện \texttt{ACTION\_UP}, \texttt{ACTION\_CANCEL}, kết thúc cuộc gọi và rời màn hình đều gửi lệnh dừng.

CaregiverApp gửi lệnh điều khiển chuẩn hóa qua Internet. RobotApp tiếp nhận, kiểm tra và chuyển lệnh thành vận tốc của hai động cơ theo quy ước lắp đặt, rồi truyền đến ESP32 qua BLE. Nhờ đó, người chăm sóc có thể điều khiển từ xa khi hai ứng dụng có kết nối Internet, còn việc giao tiếp với phần thân được thực hiện tại robot.

\subsection{Hiện thực dữ liệu dùng chung}

Các bộ lắng nghe Firebase cung cấp bản cập nhật gần thời gian thực, trong khi thao tác ghi có thể được phản ánh cục bộ trước khi đồng bộ với máy chủ \reportcite{ref:firebase_rtdb}. Vì vậy, mỗi lệnh điều khiển bắt buộc phải có timestamp; callback được gọi theo thời gian thực không có nghĩa dữ liệu nhận được luôn còn hiệu lực. Các lớp \texttt{UserProfile}, \texttt{ScheduleItem}, \texttt{FallAlertEntry}, \texttt{ReminderLogEntry} và \texttt{CallHistoryEntry} được dùng để ánh xạ dữ liệu JSON sang đối tượng trong ứng dụng.

Hồ sơ, lịch và nhật ký lời nhắc được lưu dưới các nhánh cấp gốc nhưng đều tiếp tục phân vùng theo UID. Các kênh vận hành gồm cờ chức năng, lệnh điều khiển, signaling, cảnh báo té ngã và lịch sử cuộc gọi được tập trung dưới \texttt{pair\_sessions/\{uid\}}. Security Rules áp dụng nguyên tắc từ chối mặc định và chỉ cho phép tài khoản đã xác thực truy cập nhánh có UID tương ứng \reportcite{ref:firebase_security}. Cấu trúc này ngăn dữ liệu của các tài khoản khác nhau dùng chung kênh vận hành, nhưng hai thiết bị trong một cặp vẫn đăng nhập bằng cùng tài khoản nên chưa biểu diễn được vai trò robot và người chăm sóc như hai danh tính độc lập.

\subsection{Huấn luyện và chuyển đổi mô hình AI}

Mô hình phát hiện té ngã được chuẩn bị theo một quy trình liên tục từ xử lý dữ liệu video, trích xuất đặc trưng tư thế, huấn luyện LSTM đến chuyển đổi sang định dạng phù hợp với Android. Việc thống nhất quy ước đặc trưng giữa giai đoạn huấn luyện và suy luận có ý nghĩa quyết định, bởi chỉ một sai khác về thứ tự hoặc cách chuẩn hóa tọa độ cũng làm đầu vào trên thiết bị không còn tương thích với mô hình.

\subsubsection{Xử lý dữ liệu}

Dữ liệu LE2I được cung cấp dưới dạng các video liên tục gồm hoạt động bình thường và tình huống té ngã. Quá trình tiền xử lý gồm:
\begin{itemize}
    \item \textbf{Trích xuất khung hình}: Video được phân tách thành các khung hình ở tốc độ 25 FPS để tương thích với tốc độ xử lý mục tiêu của camera robot.
    \item \textbf{Phân chia tập dữ liệu}: Tập huấn luyện, xác thực và kiểm tra được chia theo video và chủ thể, không chia ngẫu nhiên theo từng khung hình. Cách tổ chức này hạn chế rò rỉ dữ liệu khi những khung hình gần nhau của cùng một hành động xuất hiện đồng thời ở tập huấn luyện và tập kiểm tra.
    \item \textbf{Xử lý mất cân bằng lớp}: Số mẫu hoạt động bình thường thường lớn hơn số mẫu té ngã. Trọng số lớp và phương pháp lấy mẫu tăng cường được sử dụng để giảm xu hướng mô hình nghiêng về lớp NO FALL.
\end{itemize}

\subsubsection{Trích xuất đặc trưng}

Mô hình LSTM không nhận trực tiếp ảnh RGB mà học từ chuỗi đặc trưng tư thế do YOLO-Pose cung cấp. Mỗi khung hình của người mục tiêu được biểu diễn bởi 69 giá trị theo đúng quy ước sử dụng trong RobotApp:
\begin{enumerate}
    \item \textbf{Tọa độ chuẩn hóa}: Tọa độ $(x_i,y_i)$ của 17 keypoint được chuẩn hóa theo tâm hông và kích thước của vùng tư thế, giúp giảm ảnh hưởng của vị trí và khoảng cách giữa người với camera.
    \item \textbf{Đặc trưng chuyển động}: Với mỗi keypoint, hai thành phần $(\Delta x_i,\Delta y_i)$ biểu diễn độ dịch chuyển so với khung hình trước. Mỗi điểm vì vậy đóng góp bốn giá trị $(x_i,y_i,\Delta x_i,\Delta y_i)$.
    \item \textbf{Tỷ lệ hình dạng và cửa sổ thời gian}: Tỷ lệ rộng/cao của bounding box được ghép với 68 giá trị keypoint để tạo vector 69 chiều. Mười lăm khung hình liên tiếp hình thành tensor đầu vào có kích thước $[1,15,69]$ cho LSTM.
\end{enumerate}

Cấu trúc tế bào LSTM đã được giới thiệu ở Hình \ref{fig:lstm_architecture}. Trong ứng dụng, chuỗi đặc trưng tư thế được đưa vào mô hình để phân tích sự thay đổi chuyển động theo thời gian.






\subsubsection{Huấn luyện mô hình}

Mô hình LSTM được xây dựng bằng TensorFlow/Keras. Quá trình huấn luyện gồm:
\begin{itemize}
    \item Thiết lập seed cố định để tăng khả năng tái lập thí nghiệm.
    \item Sử dụng hàm mất mát Binary Cross-Entropy cho bài toán phân loại nhị phân.
    \item Áp dụng dừng sớm khi hàm mất mát trên tập xác thực không còn cải thiện, đồng thời lưu mô hình tốt nhất tại \texttt{best\_model\_lstm.h5}.
\end{itemize}

\subsubsection{Chuyển đổi và tối ưu mô hình}

Để triển khai trên điện thoại Android, mô hình sau huấn luyện được chuyển sang TensorFlow Lite nhằm giảm kích thước và tận dụng môi trường suy luận dành cho thiết bị biên.

\textbf{1. Đối với mô hình LSTM:} \\
Chương trình \path|convert_to_tflite.py| nạp \texttt{best\_model\_lstm.h5} và kiểm tra \texttt{input\_shape} trước khi chuyển đổi bằng \path|TFLiteConverter.from_keras_model|. Tập toán tử gồm TFLite built-ins và Select TF Ops. Kết quả là \texttt{best\_model\_lstm\_v2.tflite}, có thể đóng gói trực tiếp trong thư mục \texttt{assets} của RobotApp.

\textbf{2. Đối với mô hình YOLO-Pose:} \\
Mô hình PyTorch \texttt{yolo26n-pose.pt} trước hết được xuất sang ONNX với kích thước lô cố định bằng 1 và ảnh đầu vào $320\times320$. Thư viện \texttt{onnx2tf} tiếp tục chuyển mô hình ONNX sang SavedModel và TFLite. Biến thể FP16 lưu trọng số bằng số dấu phẩy động 16 bit, qua đó giảm dung lượng và băng thông bộ nhớ so với FP32; mức ảnh hưởng đến chất lượng cần được xác nhận bằng phép kiểm tra tương đương sau chuyển đổi.

\begin{table}[htbp]
\centering
\caption{Các mô hình chính sau khi huấn luyện và chuyển đổi}
\label{tab:model_artifacts}
\small
\begin{tabularx}{\textwidth}{|p{4.3cm}|p{2.5cm}|p{2.2cm}|X|}
\hline
\textbf{Tệp mô hình} & \textbf{Định dạng} & \textbf{Dung lượng} & \textbf{Mục đích} \\
\hline
best\_model\_lstm.h5 & Keras H5 & 193,8 KB & Mô hình gốc để huấn luyện và đánh giá \\
\hline
best\_model\_lstm\_v2.tflite & TFLite & 63,9 KB & Mô hình LSTM đã tối ưu cho Android \\
\hline
yolo26n-pose.pt & PyTorch & 7,88 MB & Mô hình Pose gốc \\
\hline
yolo26n-pose\_320\_fp16.tflite & TFLite FP16 & 6,28 MB & Biến thể YOLO đã lượng tử hóa FP16, tối ưu kích thước và tốc độ \\
\hline
\end{tabularx}
\end{table}

\subsubsection{Quy trình kiểm tra tương đương sau chuyển đổi}
Quy trình kiểm tra tương đương sử dụng cùng một tập dữ liệu đầu vào cho mô hình gốc và mô hình TFLite, sau đó đối chiếu bounding box, độ tin cậy, điểm khớp và vector xác suất của LSTM. Các phép so sánh nhằm phát hiện sai khác về toán tử, thứ tự đầu ra hoặc cách chuẩn hóa mà việc nạp mô hình thành công chưa phản ánh được. Đây là bước cần hoàn thiện cùng với lưu trữ dữ liệu đầu vào và kết quả so sánh để đánh giá đầy đủ ảnh hưởng của quá trình chuyển đổi.


\subsection{Hiện thực firmware phần thân robot}

\subsubsection{Khởi tạo phần cứng}

Firmware sử dụng ESP32 để giao tiếp với hai cảm biến góc AS5600 qua hai bus I2C. Cảm biến của động cơ trái nối với SDA=19/SCL=18, còn cảm biến của động cơ phải nối với SDA=23/SCL=5. Mạch công suất bên trái dùng IN1=32, IN2=33, IN3=25, EN=22; bên phải dùng IN1=26, IN2=27, IN3=14, EN=12. Hai bus I2C riêng cho phép sử dụng đồng thời hai AS5600 có cùng địa chỉ mặc định.

Hai đối tượng \texttt{BLDCMotor} được cấu hình với 7 cặp cực và liên kết với cảm biến góc cùng mạch công suất điều khiển bằng ba tín hiệu PWM. Chương trình thiết lập phương pháp điều chế \texttt{SpaceVectorPWM}, chế độ điều khiển vận tốc \texttt{velocity}, các tham số PID và bộ lọc thông thấp trước khi gọi \texttt{init()} và \texttt{initFOC()}. Trong vòng lặp chính, \texttt{loopFOC()} được gọi liên tục, không chèn thời gian chờ bằng \texttt{delay()}.

\subsubsection{Tiếp nhận lệnh qua BLE GATT}

NimBLE khởi tạo thiết bị có tên \texttt{OhmniRobot}, thiết lập công suất phát và đưa UUID dịch vụ vào dữ liệu phản hồi quét để hỗ trợ tìm kiếm. Đặc tính nhận lệnh sử dụng thuộc tính ghi không phản hồi \texttt{WRITE\_NR}. Khi kết nối được thiết lập, hàm xử lý sự kiện cập nhật trạng thái và tiếp tục phát quảng bá; khi ngắt kết nối, vận tốc đặt của hai bánh được đưa về 0. Khi nhận dữ liệu ghi, chương trình lấy chuỗi \texttt{std::string}, loại khoảng trắng ở đầu và cuối, rồi chuyển sang bộ phân tích lệnh.

Bộ phân tích tìm vị trí hai ký tự L và R, tách chuỗi con rồi chuyển sang số bằng \texttt{toFloat}. Khi dữ liệu hợp lệ, vận tốc đặt và \texttt{lastCmdTime} được cập nhật. RobotApp đã chuẩn hóa định dạng lệnh trước khi gửi; tuy nhiên, firmware vẫn cần kiểm tra toàn bộ chuỗi và dùng \texttt{isfinite} để loại giá trị không hợp lệ ngay tại tầng chấp hành.

\subsubsection{Xử lý vùng không đáp ứng và dừng khi mất lệnh}

Hàm \texttt{applyDeadzone} đưa $|v|<0,5$ về 0, nâng độ lớn của lệnh khác 0 nhưng nhỏ hơn 3 lên 3, rồi giới hạn kết quả trong $[-40,40]$. Cơ chế watchdog theo dõi thời gian kể từ lần nhận lệnh gần nhất khi BLE đang kết nối. Nếu thời gian này vượt 800 ms và vận tốc đặt khác 0, firmware ghi nhật ký và đưa vận tốc đặt của cả hai bánh về 0. Do \texttt{motor.move()} được gọi sau bước kiểm tra, lệnh dừng được áp dụng ngay trong chu kỳ điều khiển hiện tại.

\begin{lstlisting}[caption={Giả mã xử lý lệnh và điều khiển động cơ trên ESP32}]
loopFOC(left); loopFOC(right)
if connected and now - lastCommand > 800 ms:
    targetLeft = 0; targetRight = 0
leftCommand  = applyDeadzone(targetLeft)
rightCommand = applyDeadzone(targetRight)
move(leftCommand); move(rightCommand)
\end{lstlisting}

\subsection{Kết luận chương}

Chương đã trình bày quá trình hiện thực kiến trúc hệ thống thành hai ứng dụng Android, các mô hình AI và firmware ESP32. RobotApp đảm nhiệm cảm nhận, hội thoại, điều phối và kết nối với thân robot; CaregiverApp cung cấp các chức năng quản lý, giám sát và can thiệp từ xa. Firebase, WebRTC và BLE tạo thành ba kênh giao tiếp tương ứng cho dữ liệu dùng chung, truyền thông đa phương tiện và điều khiển phần cứng. Bên cạnh đó, quy trình chuyển đổi mô hình và các cơ chế dừng nhiều tầng giúp hệ thống phù hợp hơn với điều kiện vận hành trên thiết bị di động. Chương 5 tiến hành kiểm thử, đánh giá kết quả và phân tích những giới hạn của hệ thống.
